\chapter{Conclusioni}

Il progetto ha avuto come obiettivo la costruzione e l’analisi di una rete di rotte aeree tramite gli strumenti della \textit{Social Network Analysis}. A partire dal dataset iniziale, dopo una fase di preprocessing e selezione delle variabili rilevanti, è stato possibile modellare il sistema come un grafo di aeroporti e collegamenti, prima in forma diretta e successivamente anche in una versione non diretta derivata, utile per le analisi strutturali.

I risultati ottenuti hanno mostrato che la rete presenta una struttura molto ampia ma allo stesso tempo fortemente sparsa, con un basso valore di densità e una distribuzione dei gradi fortemente asimmetrica. Questo conferma la presenza di una configurazione di tipo \textit{hub-and-spoke}, in cui pochi aeroporti concentrano un numero elevato di connessioni, mentre la maggior parte dei nodi occupa posizioni più periferiche. L’analisi delle componenti ha inoltre evidenziato la presenza di una \textit{giant component} dominante, segno di un sistema globale fortemente interconnesso.

Lo studio delle misure di centralità ha permesso di mettere in luce come l’importanza di un aeroporto non possa essere descritta tramite una sola prospettiva. Alcuni nodi sono risultati particolarmente rilevanti in termini di numero di collegamenti, altri per la loro capacità di intermediazione tra porzioni differenti della rete, altri ancora per il loro elevato livello di accessibilità, per il prestigio strutturale derivante dalla qualità delle connessioni ricevute o per il fatto di essere collegati a nodi a loro volta molto centrali. In questo senso, l’uso combinato di \textit{betweenness centrality}, \textit{closeness centrality}, \textit{PageRank} ed \textit{eigenvector centrality} ha consentito di ottenere una visione più completa e articolata del ruolo svolto dai principali aeroporti nel sistema.

L’analisi strutturale ha ulteriormente arricchito la comprensione della rete. La \textit{community detection} ha mostrato che il sistema non è omogeneo, ma organizzato in gruppi di aeroporti fortemente interconnessi. L’analisi del massimo \textit{k-core} ha permesso di identificare il nucleo più coeso della rete, mentre lo studio delle \textit{ego-network} di FRA e DXB ha evidenziato in modo locale la centralità di alcuni hub particolarmente rilevanti. Anche il coefficiente di clustering ha fornito indicazioni utili sulla coesione locale, mostrando che la rete presenta aree fortemente interconnesse ma che i grandi hub globali mantengono spesso una struttura più aperta e articolata. Infine, la presenza di clique massimali di dimensione significativa conferma che all’interno della rete esistono sottostrutture ad alta coesione.

Nel complesso, il progetto ha mostrato come una rete di trasporto aereo possa essere studiata efficacemente tramite i metodi della Social Network Analysis, ricavando informazioni utili non solo dal punto di vista descrittivo, ma anche in termini di interpretazione strutturale del sistema. La rappresentazione a grafo si è dimostrata particolarmente adatta a evidenziare la distribuzione dei ruoli tra i nodi, la presenza di hub globali, l’organizzazione in comunità, la formazione di nuclei densi e il diverso contributo dei singoli aeroporti alla connettività complessiva.

Naturalmente, il lavoro presenta anche alcuni limiti. In primo luogo, il dataset considera le rotte come collegamenti tra aeroporti, ma non include informazioni dinamiche come il volume reale del traffico, la frequenza dei voli o variazioni temporali. Inoltre, il peso degli archi è stato definito come numero di compagnie distinte che operano una rotta, scelta metodologicamente sensata ma non unica possibile. Infine, alcune analisi, come la \textit{betweenness centrality}, sono state calcolate in forma approssimata per ragioni computazionali legate alla dimensione della rete.

Tra i possibili sviluppi futuri si possono considerare diverse estensioni del progetto. Una prima possibilità consiste nell’integrare informazioni geografiche o temporali, così da studiare l’evoluzione della rete o confrontare sottoreti regionali differenti. Un secondo sviluppo potrebbe riguardare l’impiego di ulteriori misure di robustezza o vulnerabilità della rete, per valutare l’impatto della rimozione di nodi particolarmente centrali. Infine, sarebbe interessante confrontare differenti definizioni del peso degli archi, ad esempio basate sul numero di voli o sul traffico passeggeri, al fine di ottenere una rappresentazione ancora più aderente al funzionamento reale del sistema aereo.